\section{Desafios práticos}

\desafio{Gere o mesmo serviço em modo JVM e nativo. Meça o tempo até a
primeira requisição e o RSS seguindo o guia de medição (não use
\texttt{docker stats}). Depois rode o teste de carga da Parte~\ref{parte:testes} por dez
minutos e compare o throughput no primeiro e no último minuto. Explique
a diferença entre os dois modos usando a seção~\ref{sec:jit}.}

\desafio{Crie uma classe com
\texttt{static final String AMBIENTE = System.getenv("AMBIENTE")} e um
endpoint que devolve o valor. Gere o executável nativo com uma variável
definida no build e rode-o com outro valor. O que o endpoint devolve?
Conserte sem tirar o campo \texttt{static}.}

\desafio{Devolva de um endpoint um objeto serializado por Jackson cuja
classe só é usada por reflexão. Encontre o erro no executável nativo e
conserte de duas formas. Qual delas você deixaria no código, e por quê?}

\desafio{Pegue um endpoint bloqueante que chama o banco e rode-o em
worker thread e com \texttt{@RunOnVirtualThread}, com 50, 500 e 5000
usuários concorrentes. Meça p99, threads de plataforma e conexões em uso
no pool. Em que ponto o gargalo deixou de ser thread e passou a ser
outra coisa?}

\desafio{Rode o serviço num container com 512~MB de limite, sem opções
de memória. Descubra o heap máximo que a JVM escolheu (\texttt{jcmd}
ou \texttt{-XX:+PrintFlagsFinal}) e compare com o RSS sob carga. Repita
com o executável nativo sem \texttt{-Xmx}. O que você passaria a
configurar explicitamente?}
